\documentclass[11pt]{article}
\usepackage[a4paper, margin=1in]{geometry}
\usepackage{amsmath, amssymb}
\usepackage{hyperref}
\usepackage{xcolor}
\usepackage{enumitem}
\usepackage{titlesec}
\usepackage{listings}
\hypersetup{colorlinks=true, linkcolor=blue!50!black, urlcolor=blue!50!black}
\titleformat{\section}{\large\bfseries\color{blue!40!black}}{\thesection.}{0.5em}{}
\titleformat{\subsection}{\normalsize\bfseries}{\thesubsection}{0.5em}{}

\definecolor{codebg}{RGB}{248,248,245}
\lstset{
  basicstyle=\ttfamily\small,
  backgroundcolor=\color{codebg},
  frame=single,
  framerule=0pt,
  rulecolor=\color{codebg},
  breaklines=true,
  language=Python,
}

\title{\bfseries Machine Learning Concepts for the WSS-from-MRI PINN Project\\[0.3em]
\large What you need to know, what you'll actually build, and what the alternatives look like}
\author{}
\date{}

\begin{document}
\maketitle

\section*{How to read this document}
This document assumes you are mathematicians comfortable with optimisation theory, linear algebra, and PDE analysis, but treating ML as a new field. The concepts are presented in the order you will encounter them when building the project. Each section ends with a \textbf{do/don't} synthesis.

\section{Neural networks as function approximators}

A feed-forward neural network with $L$ layers is a parameterised function $f_\theta: \mathbb{R}^{d_{\text{in}}} \to \mathbb{R}^{d_{\text{out}}}$ defined by
\[
\mathbf{h}_0 = \mathbf{x}, \qquad
\mathbf{h}_\ell = \phi(W_\ell \mathbf{h}_{\ell-1} + \mathbf{b}_\ell), \quad \ell = 1,\dots,L-1, \qquad
f_\theta(\mathbf{x}) = W_L \mathbf{h}_{L-1} + \mathbf{b}_L,
\]
with weights $W_\ell$, biases $\mathbf{b}_\ell$, activation $\phi$, and parameters $\theta = \{W_\ell, \mathbf{b}_\ell\}$.

\subsection*{Universal approximation}
A network with one hidden layer and a non-polynomial activation can approximate any continuous function on a compact set arbitrarily well (Cybenko 1989, Hornik 1991). Mostly a sanity check; actual performance is governed by depth, width, and \emph{trainability}.

\subsection*{Activations}
\begin{itemize}[noitemsep]
\item \textbf{ReLU} $\phi(z) = \max(0,z)$: standard, fast, but piecewise-linear (so $f_\theta$ has zero curvature almost everywhere --- bad for PINNs that need second derivatives).
\item \textbf{Tanh}, \textbf{GELU}, \textbf{Swish}: smooth alternatives. \textbf{Tanh is the default for PINNs} because it's smooth, bounded, and well-conditioned for autograd.
\item \textbf{SIREN} $\phi(z) = \sin(z)$: periodic activations; very strong at high-frequency representation but needs careful initialisation.
\end{itemize}

\subsection*{The MLP backbone you'll use}
A typical PINN backbone for vascular flow:
\begin{itemize}[noitemsep]
\item Input dimension: 3 (or 4 with time, or 5 with parameters).
\item 3--6 hidden layers, 64--256 neurons each.
\item Tanh activation.
\item Output dimension: 3 velocity components (and optionally 1 pressure component).
\end{itemize}

\textbf{Do:} start with the smallest network that fits the training data; only deepen if necessary.\\
\textbf{Don't:} use ReLU. The PDE residual involves $\partial_i u_j$ and $\partial^2_{ii} u_j$; ReLU's lack of curvature breaks the second-derivative term in the loss.

\section{Coordinate-based / implicit neural representations}

A coordinate network is just an MLP whose input is a spatial coordinate (and possibly time) and whose output is a field value. The point is that you query the network at \emph{any} point, not at fixed grid nodes. This is the right frame for vascular flow: arbitrary 3D query, no mesh.

\textit{Examples:} NeRF (3D scenes), DeepSDF (signed distance fields), and PINNs in their basic form.

\subsection*{Spectral bias}
Coordinate networks with smooth activations have an implicit \emph{low-frequency bias}: they learn smooth functions easily and high-frequency detail slowly (Rahaman et al.\ 2019, Tancik et al.\ 2020, Wang \& Perdikaris on the eigenvector bias in PINNs). The neural tangent kernel (NTK) analysis makes this precise: training reduces error in the directions of the largest NTK eigenvalues first, which correspond to low-frequency modes.

\textit{Project relevance:} the failure of vanilla PINNs near the wall is partly spectral bias --- the wall gradient is a high-frequency feature relative to the bulk flow. Mitigations matter.

\section{Random Fourier features (RFF) / positional encodings}

Replace the raw coordinate $\mathbf{r} = (x,y,z)$ with the embedded vector
\[
\boldsymbol{\phi}(\mathbf{r}) = \big[\sin(2\pi B\mathbf{r}),\;\cos(2\pi B\mathbf{r})\big] \in \mathbb{R}^{2D},
\]
where $B \in \mathbb{R}^{D\times 3}$ has entries $B_{ij}\sim \mathcal{N}(0,\sigma^2)$. The bandwidth $\sigma$ controls the highest frequency the network can represent. By Bochner's theorem, this is a finite-dimensional approximation to a Gaussian RBF kernel.

\textit{Project relevance:} this is the trick the source paper uses; it is also a candidate methodological tweak to your own work. The hyperparameters are $D$ (number of features) and $\sigma$ (bandwidth). Practical defaults: $D = 64$--$256$, $\sigma$ tuned in $\{1, 2, 5, 10\}$.

\textbf{Do:} include an RFF ablation in your paper.\\
\textbf{Don't:} apply RFF to non-spatial inputs (parameters like Re, $\sigma_v$); use raw values for those.

\section{Physics-informed neural networks (PINNs): the framework}

\subsection*{The unified loss}
A PINN trains $\mathbf{u}_\theta(\mathbf{x}, t)$ (and optionally $p_\theta$) by minimising a composite loss
\begin{equation}
\mathcal{L}(\theta) = \lambda_{\text{data}}\,\mathcal{L}_{\text{data}}(\theta) + \lambda_{\text{phys}}\,\mathcal{L}_{\text{phys}}(\theta) + \lambda_{\text{bc}}\,\mathcal{L}_{\text{bc}}(\theta),
\label{eq:PINNloss}
\end{equation}
where each term is a mean-squared error:
\begin{align*}
\mathcal{L}_{\text{data}} &= \frac{1}{N_d}\sum_{i=1}^{N_d}\|\mathbf{u}_\theta(\mathbf{x}_i, t_i) - \mathbf{u}_i^{\text{obs}}\|^2 && \text{(observation fit)}\\
\mathcal{L}_{\text{phys}} &= \frac{1}{N_r}\sum_{j=1}^{N_r}\|\mathbf{R}_{\text{NS}}(\mathbf{u}_\theta, p_\theta)(\mathbf{x}_j, t_j)\|^2 + \frac{1}{N_r}\sum_{j=1}^{N_r}|\nabla\cdot\mathbf{u}_\theta(\mathbf{x}_j,t_j)|^2 && \text{(PDE residual)}\\
\mathcal{L}_{\text{bc}} &= \frac{1}{N_b}\sum_{k=1}^{N_b}\|\mathbf{u}_\theta(\mathbf{x}_k^w, t_k) - \mathbf{0}\|^2 && \text{(no-slip)}
\end{align*}
The PDE residual $\mathbf{R}_{\text{NS}}$ is computed by autograd: the network's output is differentiated with respect to its input. In this project the three canonical terms are augmented by a pressure anchor, an inflow magnitude term, a dense-interpolant prior, and a magnitude floor --- all introduced to defeat the trivial-solution collapse (see \emph{Avoiding collapse to the trivial solution}).

\subsection*{Collocation points}
The residual is enforced at a set of \emph{collocation points} $\{\mathbf{x}_j\}$ sampled from the interior of $\Omega$. Sampling can be uniform random, Latin hypercube, importance-sampled (more density near walls / stenoses), or adaptive (residual-based).

\textbf{Do:} use $\sim 10^4$--$10^5$ collocation points per epoch; resample each epoch.\\
\textbf{Don't:} fix collocation points across epochs --- it lets the network ``cheat'' by overfitting on those.

\subsection*{Loss balancing}
The loss weights $\lambda_{\text{data}}, \lambda_{\text{phys}}, \lambda_{\text{bc}}$ are crucial. Choices range from manual tuning to learned (Wang--Teng--Perdikaris) and gradient-norm-based (Wang \emph{et al.}). Default to manual at first, then add adaptive balancing in an ablation.

\subsection*{Hard vs.\ soft constraints}
Boundary conditions can be \emph{soft} (added to the loss) or \emph{hard} (baked into the architecture). For instance, write
\[
\mathbf{u}_\theta(\mathbf{x}) = \psi(\mathbf{x}) \cdot \mathbf{N}_\theta(\mathbf{x}),
\]
where $\psi(\mathbf{x})$ is a smooth function vanishing on the wall (e.g.\ a signed-distance function $\psi(\mathbf{x}) = \mathrm{SDF}(\mathbf{x})$ inside the lumen). This guarantees no-slip exactly. Hard constraints often improve training stability and accuracy, especially in your problem because WSS recovery hinges on respecting no-slip.

\textit{Project relevance:} the hard no-slip constraint via SDF, $\mathbf{u}_\theta = \mathrm{SDF}(\mathbf{x})\cdot\mathbf{N}_\theta(\mathbf{x})$, is \textbf{adopted as the default}. WSS at the wall then follows analytically from the product rule (since $\mathrm{SDF}=0$, $\nabla\mathrm{SDF}=-\hat{\mathbf{n}}$ there): $|\boldsymbol\tau_w| = \mu(U/L)\,|\mathbf{N}_\theta^{\text{tang}}|$, needing only a forward pass.

A second hard constraint --- a \emph{vector-potential} parameterisation $\mathbf{u}=\nabla\times\mathbf{A}_\theta$, which makes $\nabla\cdot\mathbf{u}\equiv0$ by identity --- was trialled but \textbf{demoted from the default}. It requires third-order autograd for the Stokes residual (the Laplacian of a curl), which is gradient-starved and slow, and it makes WSS depend on $|\nabla\times\mathbf{A}|$ exactly where the data is sparsest (the near-wall region), so WSS could collapse \emph{independently} of the bulk velocity. The production path is therefore the SDF-hard MLP with a \emph{soft} divergence penalty; the vector potential remains available as an ablation.

\subsection*{Avoiding collapse to the trivial solution}
Because the steady Stokes operator is linear and homogeneous, the data term is the \emph{only} thing pinning the field magnitude, and a handful of sparse, noisy interior voxels cannot hold it against a physics residual that is minimised by $\mathbf{u}\to\mathbf{0}$. The adopted recipe stacks four mutually reinforcing measures:
\begin{enumerate}[noitemsep]
\item \textbf{Scale-invariant (relative) data and inflow misfit.} Normalise each fit term by the mean squared observation magnitude, $\mathcal{L}_{\text{data}} = \overline{\|\mathbf{u}_\theta-\mathbf{u}^{\text{obs}}\|^2}\,/\,(\overline{\|\mathbf{u}^{\text{obs}}\|^2}+\varepsilon)$. A zero prediction then scores $\approx1$ regardless of the tiny absolute velocity scale ($U=2\times10^{-5}$\,m/s), so $\lambda_{\text{data}}$ is comparable to the physics weight and behaves predictably across cases and voxel sizes.
\item \textbf{Inflow magnitude constraint.} Add a term matching $\mathbf{u}_\theta$ to a measured inlet velocity (sourced honestly from the near-inlet MRI voxels). An open-boundary velocity/flux condition is what makes the magnitude of a linear homogeneous field well-posed; $\mathbf{u}\equiv\mathbf{0}$ cannot satisfy a non-zero inflow.
\item \textbf{Dense interpolant prior with curriculum warm-up.} Interpolate the sparse MRI voxels into a dense interior velocity field (scattered-data radial-basis / thin-plate spline, with no-slip injected at the wall) and supervise on it. This gives a non-zero magnitude target at \emph{every} collocation point. Train data-only for a short warm-up (physics off), then ramp physics in while the interpolant weight decays linearly to zero --- the prior prevents the early slide into the trivial basin, then yields to the physics so it does not bias the final field. (Curriculum / data-first warm-up: Krishnapriyan \emph{et al.}\ 2021.)
\item \textbf{One-sided magnitude floor.} Penalise $\big[\mathrm{relu}(\text{rms}_{\text{obs}} - \text{rms}(|\mathbf{u}_\theta|))\big]^2$ over the collocation batch: zero once the field magnitude is healthy (so a correctly scaled solution is never biased), positive while collapsing. Penalising the \emph{RMS/variance} --- not the deviation from the data mean --- is essential, so spatial structure is not flattened.
\end{enumerate}
Checkpoint selection compounds these: the ``best'' model is chosen by \emph{observation misfit} (data $+$ inflow), never total loss, because the collapsed near-zero field has the \emph{lowest} total loss and would otherwise be saved as best.

\section{Inverse problems and data assimilation with PINNs}

Your problem is an inverse problem: the network is fit to noisy observations \emph{plus} the physics. The data term plays the role of a likelihood, the physics term plays the role of a regulariser, and the network parameters play the role of a high-dimensional latent state.

\subsection*{Connection to Tikhonov regularisation}
For a linear inverse problem $A\mathbf{u} = \mathbf{d}$ with noisy $\mathbf{d}$, the Tikhonov solution minimises $\|A\mathbf{u} - \mathbf{d}\|^2 + \lambda\|R\mathbf{u}\|^2$. The PINN loss \eqref{eq:PINNloss} is the same shape with $A$ being ``evaluate $\mathbf{u}_\theta$ at observation locations'' and $R$ being the PDE operator. The regularisation strength is $\lambda_{\text{phys}}/\lambda_{\text{data}}$.

\subsection*{Bayesian PINNs (B-PINNs)}
Treat $\theta$ as random. Posterior $p(\theta | \mathbf{d}) \propto p(\mathbf{d}|\theta)\,p(\theta)$ where the likelihood is the data term and the prior includes the PDE residual. Sample via Hamiltonian Monte Carlo (slow) or use a variational approximation (faster, e.g.\ MFVI, mean-field variational inference). Yang--Meng--Karniadakis 2021 is the canonical reference. Output: a distribution over $\mathbf{u}_\theta$, hence over $\boldsymbol{\tau}_w$, hence \emph{calibrated uncertainty maps}.

\textit{Project relevance:} a B-PINN extension is the strongest secondary contribution you can offer. ``This region of the wall has predicted WSS $1.2 \pm 0.4$\,Pa with credible interval [0.5, 2.1]'' is exactly what clinicians need.

\section{Training dynamics and pitfalls specific to PINNs}

\subsection*{Optimisers}
\begin{itemize}[noitemsep]
\item \textbf{Adam / AdamW}: default for the first phase of training.
\item \textbf{L-BFGS}: second-order, often used as a \emph{refinement} step after Adam. Particularly effective for PINNs because the loss landscape near a good minimum is well-conditioned.
\end{itemize}
Standard recipe: 50k--200k Adam iterations followed by L-BFGS to convergence.

\subsection*{Learning-rate schedule}
Cosine annealing or step decay. Start at $10^{-3}$, anneal to $10^{-6}$.

\subsection*{Failure modes (you will see these)}
\begin{enumerate}[noitemsep]
\item \textbf{Convergence to trivial $\mathbf{u}\equiv\mathbf{0}$.} The dominant failure mode in this project, because steady Stokes is \emph{linear and homogeneous}: $\mathbf{u}\equiv\mathbf{0},\,p\equiv\text{const}$ is an exact zero-residual solution, so physics + no-slip both reward shrinking the field. Raising $\lambda_{\text{data}}$ alone was insufficient; see \emph{Avoiding collapse to the trivial solution} below for the full anti-collapse recipe (scale-invariant data loss, inflow magnitude constraint, dense interpolant prior with curriculum warm-up, and a one-sided magnitude floor).
\item \textbf{Smooth bulk fit, wrong WSS.} The classical near-wall failure. Symptoms: low validation MSE on volume velocity but $\boldsymbol{\tau}_w$ off by a factor.
\item \textbf{Pressure indeterminacy.} Pressure is determined only up to a constant by the equations; the network may drift. Anchor pressure at one outlet point.
\item \textbf{Stiff loss landscape.} Symptom: training stalls at large loss. Causes: inadequate non-dimensionalisation; high Re; bad initialisation. Mitigations: rescale, reduce Re for development, use SIREN initialisation, gradient clipping.
\end{enumerate}

\subsection*{Computing WSS from a trained PINN}
Once $\mathbf{u}_\theta$ is trained, query at a wall point $\mathbf{x}_w$ with outward normal $\hat{\mathbf{n}}$ and compute, via autograd:
\[
\mathbf{D}_\theta(\mathbf{x}_w) = \tfrac{1}{2}\!\big(\nabla\mathbf{u}_\theta + (\nabla\mathbf{u}_\theta)^\top\big)\Big|_{\mathbf{x}_w},
\]
then $\boldsymbol{\tau}_w$ from \eqref{eq:WSS} (Document 2). \emph{The network is differentiable everywhere, by construction; this is the central PINN advantage over voxel-based finite differences.} With the hard-SDF default this reduces to the closed form $|\boldsymbol\tau_w|=\mu(U/L)\,|\mathbf{N}_\theta^{\text{tang}}|$ (no autograd through the wall); without it, the autograd Jacobian path above is used.

\section{Evaluation: more than just one MSE}

\subsection*{Field-level metrics}
\begin{itemize}[noitemsep]
\item Mean absolute error (MAE), root-mean-square error (RMSE).
\item Normalised RMSE (NRMSE): RMSE divided by the RMS of the truth. Often reported as a percentage.
\item $R^2$ coefficient of determination.
\end{itemize}

\subsection*{WSS-specific metrics}
\begin{itemize}[noitemsep]
\item Pointwise relative error in $\tau_w$ on the wall surface.
\item Mean WSS, peak WSS, area-weighted distribution metrics.
\item TAWSS and OSI (in pulsatile cases).
\end{itemize}

\subsection*{Physical-consistency metrics}
\begin{itemize}[noitemsep]
\item Cross-section flow-rate constancy: max relative deviation of $Q(z)$ along the centreline.
\item Mass-balance error at bifurcations.
\item Pressure-drop consistency: $\Delta P\cdot Q$ vs.\ volume-integrated dissipation.
\end{itemize}

\subsection*{Clinical / agreement metrics}
\begin{itemize}[noitemsep]
\item Bland--Altman analysis on TAWSS or peak WSS, network vs.\ CFD ground truth.
\item Pearson and Spearman correlation.
\item ICC for paired comparisons across geometries.
\end{itemize}

\section{Architectural alternatives you should know about}

\subsection*{DeepONet}
Learns operators $\mathcal{G}: \mathcal{F}_{\text{input}} \to \mathcal{F}_{\text{output}}$, e.g., from a boundary-condition function to a velocity field. Two networks: ``branch'' (encodes input function) and ``trunk'' (encodes evaluation point). Useful if you want a single trained model to handle many geometries.

\subsection*{Fourier Neural Operator (FNO)}
Learns operators by parameterising integral kernels in Fourier space. Resolution-invariant. Strong on Cartesian PDEs; harder on irregular geometries (Geo-FNO addresses this).

\subsection*{Graph Neural Networks (GNNs)}
Operate on unstructured meshes. PI-GNN (physics-informed GNN) approaches for WSS prediction (Wang et al.\ 2025) are competitive supervised baselines; useful as a comparison method.

\subsection*{Convolutional approaches (U-Net, super-resolution networks)}
Used in Fathi et al.\ 2020 for 4D-flow super-resolution. \emph{Supervised}, not physics-informed in the strict sense (sometimes augmented with continuity loss). Useful baseline.

\textbf{Recommendation.} Stay with a coordinate-based PINN as your primary method. Compare against (i) tri-linear interpolation of the noisy field plus finite-difference WSS, (ii) a supervised MLP (no physics), and (iii) optionally one of the operator-learning baselines. This gives you a defensible comparison story without doubling your engineering burden.

\section{Practical training hygiene}

\subsection*{Reproducibility}
\begin{itemize}[noitemsep]
\item Fix random seeds for numpy, torch, and python's \texttt{random}.
\item Log every hyperparameter via Hydra or a YAML config.
\item Track every run with W\&B or MLflow.
\item Version your synthetic-MRI generation pipeline; tiny noise differences can flip results.
\end{itemize}

\subsection*{Checkpointing and selection}
Save checkpoints every $N$ epochs; keep the one with the lowest \emph{validation physics + data loss} (don't use training loss for selection).

\subsection*{Early stopping}
Use validation loss; patience $\sim 5$--$10$\% of total epochs.

\subsection*{Hardware}
A single GPU (RTX 3090 / A100 class) is enough for 3D Stokes-regime PINNs on idealised geometry. Patient-specific aortas at higher Re may need larger memory or distributed training.

\section{What to read in the ML literature}

In addition to the must-read PINN papers in the strategy document:
\begin{itemize}[noitemsep]
\item Goodfellow, Bengio, Courville, \emph{Deep Learning} (2016) --- chapters 6--8 for the basics.
\item Tancik et al.\ 2020 (NeurIPS) --- ``Fourier features let networks learn high-frequency functions'' --- read this carefully; it's directly applicable.
\item Wang, Wang, Perdikaris 2021 (CMAME) --- ``On the eigenvector bias of Fourier feature networks: from regression to solving multi-scale PDEs with PINNs'' --- the bridge to your problem.
\item Karniadakis et al.\ 2021 (Nat.\ Rev.\ Phys.) --- ``Physics-informed machine learning'' --- canonical review.
\item Cai et al.\ 2022 ``Physics-informed neural networks (PINNs) for fluid mechanics'' (Acta Mech. Sin.) --- domain-specific review.
\item Krishnapriyan et al.\ 2021 (NeurIPS) --- ``Characterizing possible failure modes in physics-informed neural networks'' --- motivates the curriculum / data-first warm-up used in \S\ref{sec:anticollapse}.
\end{itemize}

\section*{Synthesis: what your tech stack should look like}

\begin{itemize}[noitemsep]
\item \textbf{Framework:} PyTorch (autograd is essential).
\item \textbf{PINN scaffolding:} DeepXDE for the first prototype, then your own PyTorch implementation.
\item \textbf{CFD ground truth:} FEniCSx for idealised geometries; SimVascular for VMR cases.
\item \textbf{Synthetic-MRI pipeline:} custom Python; voxelisation via PyVista or scipy; noise via numpy.
\item \textbf{Visualisation:} PyVista (3D), matplotlib (2D), seaborn (statistics).
\item \textbf{Experiment tracking:} W\&B + Hydra config.
\item \textbf{Compute:} 1 GPU sufficient for idealised; cluster access useful for the systematic sweep.
\end{itemize}

\textbf{Final do/don't.}\\
\textbf{Do:} non-dimensionalise; use tanh; sample collocation points fresh each epoch; track every hyperparameter; ablate (RFF, hard constraints, Bayesian).\\
\textbf{Don't:} use ReLU; trust a low MSE without checking WSS; skip pressure anchoring; use one random seed.

\end{document}
